\documentclass[12pt]{article}
\usepackage[utf8]{inputenc}
\usepackage[T1]{fontenc}
\usepackage[margin=1in]{geometry}
\usepackage{setspace}
\usepackage{xcolor}
\usepackage[colorlinks=true, urlcolor=blue!60!black, linkcolor=blue!60!black]{hyperref}
\setlength{\parindent}{2em}
\setlength{\parskip}{1.25em}
\renewcommand{\baselinestretch}{1.0}
\begin{document}
\begin{center}
{\LARGE \textit{Marbury v. Madison}, 5 U.S. (1 Cranch) 137 (1803)\par}
\vspace{0.8em}
{\large Case Brief\par}
\end{center}
\vspace{1em}

\section{Memory Jogger}

Federal courts have the power and duty of judicial review to refuse to enforce acts of Congress that conflict with the Constitution, and Congress cannot expand the Supreme Court's original jurisdiction beyond the categories enumerated in Article III, Section 2.

\section{Detailed Case Facts}

In the bitterly contested presidential election of 1800, Democratic-Republican Thomas Jefferson defeated Federalist incumbent John Adams, and Democratic-Republicans won control of Congress (\href{https://supreme.justia.com/cases/federal/us/5/137/}{Justia}; \href{https://www.britannica.com/event/Marbury-v-Madison}{Britannica}). During the lame-duck period before Jefferson's inauguration on March 4, 1801, the outgoing Federalist Congress passed both the Judiciary Act of 1801 (reorganizing the federal circuit courts) and the Organic Act of the District of Columbia of February 27, 1801, ch. 15, \S~11, 2 Stat. 103, 107, which authorized the President to appoint ``such number of discreet persons to be justices of the peace as the President of the United States shall, from time to time, think expedient, to continue in office for five years'' (\href{https://supreme.justia.com/cases/federal/us/5/137/}{5 U.S. at 154}; \href{https://scholarship.law.umn.edu/cgi/viewcontent.cgi?article=2064&context=concomm}{Constitutional Commentary}).

On March 2--3, 1801, President Adams nominated forty-two justices of the peace for the District of Columbia, including William Marbury (for Washington County) and three future co-petitioners: Dennis Ramsay, Robert Townsend Hooe, and William Harper (\href{https://supreme.justia.com/cases/federal/us/5/137/}{5 U.S. at 137--38, 154--55}; \href{https://fedsoc.org/case/marbury-v-madison}{Federalist Society}). The Senate confirmed the nominations on March 3, 1801, President Adams signed the commissions, and outgoing Secretary of State John Marshall---who had already been appointed Chief Justice of the United States and was temporarily holding both offices---affixed the Great Seal of the United States to the commissions (\href{https://supreme.justia.com/cases/federal/us/5/137/}{5 U.S. at 155}; \href{https://en.wikipedia.org/wiki/Marbury_v._Madison}{Wikipedia}).

In the rush of the final hours before Adams left office, Marshall's brother, James Markham Marshall, delivered most of the commissions, but several---including Marbury's---remained undelivered in the State Department office when Jefferson took office on March 4, 1801 (\href{https://en.wikipedia.org/wiki/Marbury_v._Madison}{Wikipedia}). Treating the undelivered commissions as incomplete and void, President Jefferson instructed his new Secretary of State, James Madison (and initially acting Secretary of State Levi Lincoln), to withhold them (\href{https://supreme.justia.com/cases/federal/us/5/137/}{Justia}; \href{https://constitutioncenter.org/the-constitution/supreme-court-case-library/marbury-v-madison}{National Constitution Center}).

\section{Procedural History}

At the Supreme Court's December 1801 Term, William Marbury, Dennis Ramsay, Robert Townsend Hooe, and William Harper, represented by former Attorney General Charles Lee, filed an original petition directly in the U.S. Supreme Court seeking a rule to show cause why a writ of mandamus should not issue compelling Secretary of State James Madison to deliver their commissions (\href{https://supreme.justia.com/cases/federal/us/5/137/}{5 U.S. at 137--38}; \href{https://www.courtlistener.com/opinion/84759/marbury-v-madison/}{CourtListener}). They invoked Section 13 of the Judiciary Act of 1789, 1 Stat. 73, 80--81, which authorized the Supreme Court to issue writs of mandamus ``to any courts appointed, or persons holding office, under the authority of the United States'' (\href{https://supreme.justia.com/cases/federal/us/5/137/}{5 U.S. at 173}; \href{https://www.law.cornell.edu/constitution-conan/article-3/section-1/marbury-v-madison-and-judicial-review}{Constitution Annotated}).

The Supreme Court issued the rule to show cause in December 1801 (\href{https://supreme.justia.com/cases/federal/us/5/137/}{5 U.S. at 138, 153}). Before the Court could hear the case, however, the Democratic-Republican Congress repealed the Judiciary Act of 1801 and enacted the Judiciary Act of 1802, which abolished the Court's June and December 1802 terms and postponed its next sitting until February 1803 (\href{https://supreme.justia.com/cases/federal/us/5/137/}{Justia}; \href{https://en.wikipedia.org/wiki/Marbury_v._Madison}{Wikipedia}). When arguments finally proceeded on February 11, 1803, neither Madison nor counsel for the executive appeared to show cause; instead, State Department clerks Jacob Wagner and Daniel Brent, along with Attorney General Levi Lincoln, were called as witnesses regarding the commissions' existence (\href{https://supreme.justia.com/cases/federal/us/5/137/}{5 U.S. at 139}; \href{https://www.oyez.org/cases/1789-1850/5us137}{Oyez}). The Court issued its decision on February 24, 1803 (\href{https://www.oyez.org/cases/1789-1850/5us137}{Oyez}).

\section{Judicial Votes}

The Supreme Court decided \textit{Marbury v. Madison} by a unanimous \textbf{4--0} vote of the four participating justices (the Court then consisted of six seats) (\href{https://constitutioncenter.org/the-constitution/supreme-court-case-library/marbury-v-madison}{National Constitution Center}; \href{https://www.annenbergclassroom.org/resource/the-pursuit-of-justice/pursuit-justice-chapter-1-rise-judicial-review/}{Annenberg Classroom}).

\begin{itemize}
\item \textbf{Majority (4):} Chief Justice John Marshall (author of the Opinion of the Court), joined by Justices William Paterson, Samuel Chase, and Bushrod Washington (\href{https://constitutioncenter.org/the-constitution/supreme-court-case-library/marbury-v-madison}{National Constitution Center}).
\item \textbf{Concurring / Dissenting (0):} None.
\item \textbf{Not Participating (2):} Justices William Cushing and Alfred Moore were absent due to illness and did not participate (\href{https://www.annenbergclassroom.org/resource/the-pursuit-of-justice/pursuit-justice-chapter-1-rise-judicial-review/}{Annenberg Classroom}; \href{https://en.wikipedia.org/wiki/Marbury_v._Madison}{Wikipedia}).
\end{itemize}

\section{Holding}

\textbf{Judgment:} Rule to show cause discharged; petition for a writ of mandamus dismissed for lack of original jurisdiction (\href{https://supreme.justia.com/cases/federal/us/5/137/}{5 U.S. at 180}).

\textbf{Rule Announced:} (1) An act of Congress repugnant to the U.S. Constitution is void, and federal courts have the constitutional authority and duty---judicial review---to determine whether a statute conflicts with the Constitution and to refuse to apply an unconstitutional statute in a case before them (\href{https://supreme.justia.com/cases/federal/us/5/137/}{5 U.S. at 177--80}). (2) Article III, Section 2 sets an exclusive ceiling on the Supreme Court's original jurisdiction; Congress may not enlarge that original jurisdiction by statute, rendering Section 13 of the Judiciary Act of 1789 unconstitutional insofar as it purported to grant the Supreme Court original jurisdiction to issue writs of mandamus to federal executive officers (\href{https://supreme.justia.com/cases/federal/us/5/137/}{5 U.S. at 173--76}).

\textbf{Holding vs. Dicta:} Because the Court ultimately concluded that it lacked original jurisdiction to entertain Marbury's petition, its jurisdictional ruling---that Section 13 of the Judiciary Act of 1789 unconstitutionally expanded the Court's Article III original jurisdiction and must be disregarded under the doctrine of judicial review---is the strict \textbf{holding} (\href{https://supreme.justia.com/cases/federal/us/5/137/}{5 U.S. at 173--80}). By contrast, Chief Justice Marshall's lengthy threshold determinations on Questions 1 and 2---that Marbury's appointment vested upon signature and sealing, that Madison acted unlawfully in withholding the commission, and that a writ of mandamus is an appropriate judicial remedy against executive officers performing non-discretionary ministerial duties (though not discretionary political duties)---are classically viewed by scholars and critics (starting with President Jefferson) as \textbf{dicta}, because a court without jurisdiction ordinarily cannot adjudicate the merits (\href{https://scholarship.law.duke.edu/dlj/vol18/iss1/1/}{Van Alstyne, 1969 Duke L.J. 1}; \href{https://scholarship.law.missouri.edu/cgi/viewcontent.cgi?article=4238&context=mlr}{Missouri Law Review}). Nonetheless, those principles became foundational to American administrative law and the political question doctrine (\href{https://scholarship.law.columbia.edu/context/faculty_scholarship/article/1878/viewcontent/37_J._Marshall_L._Rev._481.pdf}{Merrill, 37 J. Marshall L. Rev. 481 (2004)}).

\section{Analysis of Opinions}

\subsection{Majority Opinion (Chief Justice Marshall)}

Chief Justice Marshall organized the unanimous opinion around three sequential questions, strategically inverting the usual rule that a court addresses jurisdiction before the merits (\href{https://supreme.justia.com/cases/federal/us/5/137/}{5 U.S. at 154}).

\textbf{1. Did Marbury have a vested legal right to the commission?} Yes. Analyzing Article II, Sections 2 and 3, together with the 1789 statute governing the Department of State, Marshall distinguished three steps: nomination, appointment, and commission (\href{https://supreme.justia.com/cases/federal/us/5/137/}{5 U.S. at 155--56}). Once the President signed the commission---the last voluntary act required of him---the appointment was complete, and the Secretary of State's affixing of the Great Seal provided conclusive evidence of it (\href{https://supreme.justia.com/cases/federal/us/5/137/}{5 U.S. at 157--59}). Rejecting the analogy to a private deed requiring physical delivery, Marshall reasoned that transmission was merely a matter of convenience, not a legal prerequisite to title (\href{https://supreme.justia.com/cases/federal/us/5/137/}{5 U.S. at 159--61}). Because the Organic Act of 1801 fixed a five-year term not removable at executive will, Marbury's right had vested, and withholding the commission violated that vested legal right (\href{https://supreme.justia.com/cases/federal/us/5/137/}{5 U.S. at 162}).

\textbf{2. Do the laws afford Marbury a remedy?} Yes. Invoking Blackstone's maxim that where there is a legal right there is a legal remedy, Marshall declared that the United States is ``a government of laws, and not of men'' (\href{https://supreme.justia.com/cases/federal/us/5/137/}{5 U.S. at 163}). He then drew a landmark line between \textit{discretionary political acts} and \textit{ministerial legal duties}: where an executive officer exercises constitutional or legal discretion as the President's political agent, the officer's acts are ``only politically examinable,'' and ``[q]uestions, in their nature political \dots can never be made in this court'' (\href{https://supreme.justia.com/cases/federal/us/5/137/}{5 U.S. at 166, 170}). Conversely, where the law assigns a specific ministerial duty on which individual vested rights depend---such as recording or delivering a completed commission---the officer is ``amenable to the laws'' and subject to judicial compulsion (\href{https://supreme.justia.com/cases/federal/us/5/137/}{5 U.S. at 166, 170--71}).

\textbf{3. Can the Supreme Court issue the writ of mandamus?} No. Although mandamus was the proper writ in principle (since an action of detinue could not value a public office), the Supreme Court lacked constitutional jurisdiction to issue it as an original proceeding (\href{https://supreme.justia.com/cases/federal/us/5/137/}{5 U.S. at 173--76}). Marshall read Section 13 of the Judiciary Act of 1789 as conferring original jurisdiction on the Supreme Court to issue writs of mandamus to public officers (\href{https://supreme.justia.com/cases/federal/us/5/137/}{5 U.S. at 173}). Yet Article III, Section 2 grants the Court original jurisdiction only in ``all cases affecting ambassadors, other public ministers and consuls, and those in which a state shall be a party,'' and appellate jurisdiction ``[i]n all other cases'' (\href{https://supreme.justia.com/cases/federal/us/5/137/}{5 U.S. at 174}). Rejecting the argument that Congress could add to original jurisdiction, Marshall reasoned that such a reading would render Article III's careful distribution ``mere surplusage'' and that affirmative words must here carry a negative, exclusive meaning (\href{https://supreme.justia.com/cases/federal/us/5/137/}{5 U.S. at 174--75}). Because issuing mandamus to an executive officer creates a new cause rather than revising lower-court proceedings, it was an exercise of original, not appellate, jurisdiction (\href{https://supreme.justia.com/cases/federal/us/5/137/}{5 U.S. at 175--76}).

Faced with a direct conflict between Section 13 and Article III, Marshall delivered his celebrated defense of \textbf{judicial review}. He grounded the doctrine in both structural theory and constitutional text: (a) a written constitution of limited powers is ``the fundamental and paramount law of the nation,'' meaning any repugnant statute is void (\href{https://supreme.justia.com/cases/federal/us/5/137/}{5 U.S. at 176--77}); (b) ``[i]t is emphatically the province and duty of the Judicial Department to say what the law is,'' which requires resolving conflicts between superior and inferior rules (\href{https://supreme.justia.com/cases/federal/us/5/137/}{5 U.S. at 177--78}); (c) Article III extends judicial power to ``all cases arising under the Constitution,'' and specific provisions (such as the Treason Clause's two-witness rule and bans on export taxes, bills of attainder, and ex post facto laws) are addressed directly to courts (\href{https://supreme.justia.com/cases/federal/us/5/137/}{5 U.S. at 178--79}); (d) the judicial oath under Article VI would be ``worse than solemn mockery'' if judges had to enforce unconstitutional laws (\href{https://supreme.justia.com/cases/federal/us/5/137/}{5 U.S. at 180}); and (e) the Supremacy Clause elevates only federal laws made ``in pursuance of the Constitution'' (\href{https://supreme.justia.com/cases/federal/us/5/137/}{5 U.S. at 180}). Thus, Section 13 was void, and the rule was discharged (\href{https://supreme.justia.com/cases/federal/us/5/137/}{5 U.S. at 180}).

\subsection{Concurring Opinions}

None. All four participating justices joined Chief Justice Marshall's opinion in full.

\subsection{Dissenting Opinions}

None.

\section{Examples: Future Applications}

\begin{enumerate}
\item \textbf{Same-Side Application (Congressional Expansion of Supreme Court Original Jurisdiction):} Congress enacts a statute granting the U.S. Supreme Court exclusive original jurisdiction over all constitutional challenges to federal election laws brought by private voters against the Attorney General. \textit{Why it falls on the same side:} Because the suit involves neither foreign diplomats nor a state as a party, the statute directly violates \textit{Marbury}'s core holding that Article III, Section 2 sets an exclusive ceiling on the Supreme Court's original jurisdiction, and the Court must invalidate the jurisdictional grant and dismiss the original petition.

\item \textbf{Same-Side Application (Invalidation of Unconstitutional Federal Statute):} Congress passes a statute permitting criminal conviction for treason based on the uncorroborated testimony of a single witness, and a defendant convicted under the statute appeals to the federal courts. \textit{Why it falls on the same side:} Using the exact example Chief Justice Marshall posed in \textit{Marbury} (\href{https://supreme.justia.com/cases/federal/us/5/137/}{5 U.S. at 179}), the statute directly conflicts with Article III, Section 3's two-witness requirement; under judicial review, the court must apply the paramount Constitution and vacate the conviction.

\item \textbf{Opposite-Side Application (Congressional Regulation of Appellate Jurisdiction and Lower Courts):} Congress passes a statute vesting federal district courts (rather than the Supreme Court in the first instance) with original jurisdiction to issue writs of mandamus compelling cabinet secretaries to perform nondiscretionary statutory duties, with appellate review in the courts of appeals and the Supreme Court. \textit{Why it falls on the opposite side:} Unlike Section 13 in \textit{Marbury}, this statute does not enlarge the Supreme Court's original jurisdiction; Congress has broad Article III power to vest original federal-question jurisdiction in inferior federal courts and to channel Supreme Court review through its appellate jurisdiction.

\item \textbf{Opposite-Side Application (Challenge to Discretionary Executive Foreign Policy):} A private citizen sues the Secretary of State for a writ of mandamus ordering the United States to recognize a breakaway foreign republic and negotiate a trade treaty with it. \textit{Why it falls on the opposite side:} Under \textit{Marbury}'s own political-discretion carveout (\href{https://supreme.justia.com/cases/federal/us/5/137/}{5 U.S. at 166, 170}), diplomatic recognition and treaty negotiation are political powers entrusted exclusively to executive discretion rather than ministerial duties creating vested individual rights, so courts will refuse to issue mandamus.

\item \textbf{Fence-Sitter Application (Departmentalism vs. Judicial Supremacy Over Coordinate Branches):} The Supreme Court invalidates a federal statute in a dispute between private parties, and the President---while complying with the judgment in that specific case---publicly announces that the Executive Branch retains independent authority to interpret the Constitution in non-litigated contexts and will continue enforcing an analogous provision against non-parties unless separately enjoined. \textit{Why it is a fence-sitter:} \textit{Marbury} clearly establishes \textit{judicial review} (a court's duty to apply the Constitution in deciding a case), but scholars vigorously debate whether \textit{Marbury} entails \textit{judicial supremacy} (making the Court's constitutional interpretations binding on coordinate branches outside the judgment) or permits \textit{departmentalism}, under which each branch interprets the Constitution within its own sphere (\href{https://scholarship.law.umn.edu/cgi/viewcontent.cgi?article=1616&context=concomm}{Constitutional Commentary}; \href{https://www.cambridge.org/core/journals/perspectives-on-politics/article/reconsidering-judicial-supremacy-from-the-countermajoritarian-difficulty-to-constitutional-transformations/441C07D9E0B84500BEC580C7C179341F}{Perspectives on Politics}).
\end{enumerate}

\section{Critique}

Although \textit{Marbury v. Madison} is universally celebrated as the cornerstone of American constitutionalism, both its legal reasoning and its institutional legacy have drawn penetrating critiques from across the scholarly spectrum.

In his classic doctrinal critique, \textit{A Critical Guide to Marbury v. Madison}, William W. Van Alstyne demonstrated that nearly every step of Marshall's opinion was contestable (\href{https://scholarship.law.duke.edu/dlj/vol18/iss1/1/}{Van Alstyne, 1969 Duke L.J. 1}). Marshall arguably should have recused himself because his own failure as Secretary of State to deliver the commissions caused the dispute; he inverted the normal order of decision by lecturing President Jefferson on the merits before dismissing for lack of jurisdiction; he could have read Section 13 of the Judiciary Act of 1789 narrowly as authorizing mandamus only as an ancillary remedy in cases where the Court already had jurisdiction; and his textual arguments---such as the judicial oath and ``in pursuance thereof''---beg the question of \textit{who} has final interpretive authority, since legislators and the President also swear to uphold the Constitution (\href{https://scholarship.law.duke.edu/dlj/vol18/iss1/1/}{Van Alstyne}).

From an \textbf{originalist and textualist} perspective, scholars such as Saikrishna B. Prakash, John C. Yoo, and William Michael Treanor defend the historical legitimacy of judicial review itself, showing that Founding-era constitutional text (Article III's ``judicial Power'' and the Supremacy Clause), \textit{Federalist} No. 78, and numerous pre-1803 state and federal decisions already recognized judicial review before \textit{Marbury} (\href{https://chicagounbound.uchicago.edu/uclrev/vol70/iss3/3/}{Prakash \& Yoo, 70 U. Chi. L. Rev. 887 (2003)}; \href{https://scholarship.law.georgetown.edu/context/facpub/article/2040/viewcontent/58StanLRev455.pdf}{Treanor, 58 Stan. L. Rev. 455 (2005)}). Many originalists, however, emphasize that \textit{Marbury} justified judicial review precisely because judges are bound to enforce the fixed, written meaning of the Constitution as the people's agents, rather than inventing evolving norms (\href{https://www.law.upenn.edu/live/files/47-coan158upalrev10252010pdf}{Coan, 158 U. Pa. L. Rev. 1025 (2010)}).

From a \textbf{progressive, populist, and institutional} perspective, scholars have focused on what Alexander Bickel famously termed the ``countermajoritarian difficulty''---the tension between unelected, life-tenured judges invalidating laws enacted by democratically accountable majorities (\href{http://www.princeton.edu/~kewhitt/judicial_review}{Whittington, Theories of Judicial Review}; \href{https://ilr.law.uiowa.edu/sites/ilr.law.uiowa.edu/files/2025-11/A3_Dodson_1.pdf}{Iowa Law Review}). Critics like Sylvia Snowiss and proponents of ``popular constitutionalism'' argue that modern judicial supremacy---treating the Supreme Court as the exclusive, final expositor of the Constitution for all branches---far exceeds the narrower, defensive judicial review Marshall articulated in 1803 (\href{https://scholarship.law.umn.edu/cgi/viewcontent.cgi?article=2076&context=concomm}{Snowiss, 20 Const. Comment. 231 (2003)}; \href{https://constitutioncenter.org/blog/marbury-v-madison-and-the-independent-supreme-court}{National Constitution Center}). At the same time, many modern defenders view \textit{Marbury}'s assertion of independent judicial review as indispensable for protecting minority rights, policing separation of powers, and preserving the rule of law against majoritarian or executive overreach (\href{https://www.archives.gov/milestone-documents/marbury-v-madison}{National Archives}).

Ultimately, \textit{Marbury} embodies a perpetual tension between \textbf{competing constitutional values}: constitutional supremacy and the protection of individual legal rights through an independent judiciary versus democratic self-rule, coordinate-branch equality, and judicial restraint.

\section{Key Quotations}

\begin{quote}
``The very essence of civil liberty certainly consists in the right of every individual to claim the protection of the laws whenever he receives an injury.''
\end{quote}
\noindent --- Chief Justice John Marshall, \textit{Marbury v. Madison}, 5 U.S. (1 Cranch) 137, 163 (1803) (\href{https://supreme.justia.com/cases/federal/us/5/137/}{Justia}; \href{https://www.courtlistener.com/opinion/84759/marbury-v-madison/}{CourtListener}).

\begin{quote}
``The Government of the United States has been emphatically termed a government of laws, and not of men. It will certainly cease to deserve this high appellation if the laws furnish no remedy for the violation of a vested legal right.''
\end{quote}
\noindent --- Chief Justice John Marshall, \textit{Marbury v. Madison}, 5 U.S. (1 Cranch) 137, 163 (1803) (\href{https://supreme.justia.com/cases/federal/us/5/137/}{Justia}; \href{https://www.courtlistener.com/opinion/84759/marbury-v-madison/}{CourtListener}).

\begin{quote}
``The province of the Court is solely to decide on the rights of individuals, not to enquire how the Executive or Executive officers perform duties in which they have a discretion. Questions, in their nature political or which are, by the Constitution and laws, submitted to the Executive, can never be made in this court.''
\end{quote}
\noindent --- Chief Justice John Marshall, \textit{Marbury v. Madison}, 5 U.S. (1 Cranch) 137, 170 (1803) (\href{https://supreme.justia.com/cases/federal/us/5/137/}{Justia}; \href{https://www.courtlistener.com/opinion/84759/marbury-v-madison/}{CourtListener}).

\begin{quote}
``The Constitution is either a superior, paramount law, unchangeable by ordinary means, or it is on a level with ordinary legislative acts, and, like other acts, is alterable when the legislature shall please to alter it.''
\end{quote}
\noindent --- Chief Justice John Marshall, \textit{Marbury v. Madison}, 5 U.S. (1 Cranch) 137, 177 (1803) (\href{https://supreme.justia.com/cases/federal/us/5/137/}{Justia}; \href{https://www.courtlistener.com/opinion/84759/marbury-v-madison/}{CourtListener}).

\begin{quote}
``It is emphatically the province and duty of the Judicial Department to say what the law is. Those who apply the rule to particular cases must, of necessity, expound and interpret that rule. If two laws conflict with each other, the Courts must decide on the operation of each.''
\end{quote}
\noindent --- Chief Justice John Marshall, \textit{Marbury v. Madison}, 5 U.S. (1 Cranch) 137, 177 (1803) (\href{https://supreme.justia.com/cases/federal/us/5/137/}{Justia}; \href{https://www.courtlistener.com/opinion/84759/marbury-v-madison/}{CourtListener}).

\end{document}
